We test free relaxation on the fixed $K=8$ mesh, with $T=0.5$
and $\dt=1/16,1/32,1/64$. All physical parameters remain one.
Unlike the manufactured-solution test, the applied field and all
compensating sources are set to zero:
$\bs H_e=\bs f=\bs g=\bs r=0$.
We retain the homogeneous boundary traces specified in
\Cref{subsec:experiment-setup}.
The initial shapes are
\[
\bs u_0=C_u^{-1}\nabla B\times a,\qquad
\bs\omega_0=C_\omega^{-1}B b,\qquad
\bs m_0=C_m^{-1}S c.
\]
The initial velocity is projected onto the discrete solenoidal space;
angular velocity and magnetization are interpolated into their
respective spaces. The initial potential is recomputed with zero
boundary trace from
\[
(\nabla\varphi_h^0,\nabla\psi_h)
=-(\bs m_h^0,\nabla\psi_h)
\qquad\text{for all admissible }\psi_h,
\]
rather than prescribed from the manufactured magnetic potential.
The auxiliary variables satisfy their discrete projection relations.
Thus the internal magnetic field is not artificially suppressed,
and the subsequent solution is not constrained to follow the
manufactured temporal profile.

At every time level we integrate the discrete energy
\begin{equation}\label{eq:energy-test-observable}
E_h^n=\frac12\left(
\rho\|\bs u_h^n\|^2+
\rho\kappa\|\bs\omega_h^n\|^2+
\|\bs m_h^n\|^2+
\mu_0\|\nabla\varphi_h^n\|^2\right).
\end{equation}
The norms in \eqref{eq:energy-test-observable} are spatial $L^2$ norms.
Energy and dissipation are evaluated in the original finite-element
spaces using degree-12 quadrature, not from the visualization fields.
For the first-order scheme, we check the balance defect
\begin{equation}\label{eq:energy-test-first-balance}
R_h^n=E_h^n-E_h^{n-1}+\dt D_h^n+N_h^n,
\qquad
N_h^n=\tfrac12\|\mathcal X_h^n-\mathcal X_h^{n-1}\|_{\mathtt D}^2,
\end{equation}
where $D_h^n$ is the discrete viscous, spin, magnetic-diffusion,
and relaxation dissipation corresponding to $\mathtt A$.
In particular, the weak magnetic diffusion terms are evaluated as
$\sigma\|\bs k_h^n\|^2$ and
$\sigma\mu_0\|\vartheta_h^n\|^2$.
For the second-order discretization, the balance defect is
\begin{equation}\label{eq:energy-test-second-balance}
R_h^n=E_h^n-E_h^{n-1}+\dt D_h^{n-1/2},
\end{equation}
where $D_h^{n-1/2}$ is evaluated at the midpoint averages of the
primary fields and the endpoint auxiliaries
$\bs k_h,\vartheta_h$.
There is no first-order numerical-dissipation term $N_h^n$.
The stage fields $\bs z_h,\varpi_h$ are not averaged a second time.

For both schemes, the reported balance diagnostic is
\[
\epsilon_{\rm bal}^n=
\frac{|R_h^n|}{\max\{E_h^n,E_h^{n-1},\dt D_h^n+N_h^n,
10^{-14}E_h^0\}},
\]
with the midpoint dissipation and $N_h^n=0$ used for the
second-order scheme.

Both schemes retain the spatial approximation orders specified in
\Cref{subsec:experiment-setup}. Thus their discrete initial states differ,
whereas each scheme uses the same initial state for all three time steps.

All energy runs use 48 MPI processes with one thread each and the
same five-block solver settings as the convergence experiments.
The first-order runs require no nonlinear startup.

For the second-order energy test, the initial Broyden solve uses relative tolerance
$10^{-10}$ with at most 20 iterations, and the accepted midpoint
equations are independently reassembled to check their residual.
Subsequent time steps use the prescribed extrapolation.
